\documentclass[]{report}
\usepackage{amsmath}
\usepackage[spanish]{babel}
\usepackage[utf8]{inputenc}
\usepackage{graphicx}
\usepackage{listings}
\usepackage{xcolor}
\usepackage{booktabs}
\usepackage{float}
\usepackage[top=2cm, bottom=2cm, left=3cm, right=3cm]{geometry}

\lstset{
  language=Python,
  basicstyle=\ttfamily\footnotesize,
  keywordstyle=\color{blue},
  commentstyle=\color{gray},
  stringstyle=\color{orange!80!black},
  backgroundcolor=\color{gray!8},
  frame=single,
  breaklines=true,
  showstringspaces=false,
  numbers=left,
  numberstyle=\tiny\color{gray},
  numbersep=5pt
}

\title{Introducción al Modelado Continuo \\[0.5cm]
    Trabajo Práctico 2 \\[0.5cm]
    Identificación de Número Telefónico}

\author{
    \textbf{Grupo: 11} \\[1cm]
    \begin{tabular}{lll}
        \textbf{Apellido} & \textbf{Nombre} & \textbf{Libreta} \\ \hline
        Wilders Azara & Santiago & 350/19 \\
        Peralta       & Martín   & 428/22 \\
        Vernay        & Jerónimo & 313/24 \\
    \end{tabular}
}

\date{\today}
\setlength{\parindent}{0pt}
\setlength{\parskip}{1ex plus .5ex minus .2ex}

\begin{document}
\maketitle

% ─────────────────────────────────────────────────────────────
\section*{Introducción}

El sistema DTMF (\textit{Dual-Tone Multi-Frequency}) codifica cada dígito telefónico
como la superposición de dos sinusoides puras: una frecuencia de fila y una de columna,
según la tabla:

\begin{table}[H]
\centering
\begin{tabular}{c|ccc}
\toprule
         & \textbf{1209 Hz} & \textbf{1336 Hz} & \textbf{1477 Hz} \\
\midrule
\textbf{697 Hz} & 1 & 2 & 3 \\
\textbf{770 Hz} & 4 & 5 & 6 \\
\textbf{852 Hz} & 7 & 8 & 9 \\
\textbf{941 Hz} & * & 0 & \# \\
\bottomrule
\end{tabular}
\caption{Tabla de frecuencias DTMF.}
\end{table}

El objetivo es identificar los dígitos de los archivos \texttt{tlfn-a.wav} y
\texttt{tlfn-b.wav} utilizando la transformada discreta de Fourier (DFT) calculada
mediante el algoritmo FFT.

% ─────────────────────────────────────────────────────────────
\section*{Herramienta matemática: la DFT}

Dado que trabajamos con señales digitales, disponemos de una secuencia finita de
muestras $f[0], f[1], \ldots, f[N-1]$ obtenidas a una tasa de muestreo $f_s$.
La \textbf{transformada de Fourier discreta} asocia a $f$ otra función $\hat{f}$
definida como (ec.~14.2.1 de las notas):

\begin{equation}
\hat{f}[k] = \sum_{j=0}^{N-1} f[j]\, e^{-2\pi i jk/N}, \qquad k = 0, 1, \ldots, N-1.
\label{eq:dft}
\end{equation}

Cada coeficiente $\hat{f}[k]$ mide la contribución de la frecuencia discreta
$\xi_k = k\,f_s / N$ a la señal. La resolución frecuencial es por lo tanto
\begin{equation}
\Delta\xi = \frac{f_s}{N},
\label{eq:resolucion}
\end{equation}
de modo que segmentos más largos (mayor $N$) producen una separación más fina
entre las frecuencias accesibles.

La DFT puede recuperarse a partir de sus coeficientes mediante la fórmula de
inversión (ec.~14.4.1):
\begin{equation}
f[k] = \frac{1}{N}\sum_{j=0}^{N-1} \hat{f}[j]\, e^{2\pi ijk/N}.
\end{equation}

Computacionalmente, la evaluación directa de \eqref{eq:dft} requiere $O(N^2)$
operaciones. La \textbf{transformada rápida de Fourier (FFT)}, basada en la
factorización recursiva de la matriz $A(\omega)$ (Sección~14.8), reduce este
costo a $O(N \log N)$, lo que la hace practicable para señales de audio de
decenas de miles de muestras.

Un aspecto a tener en cuenta es el \textit{aliasing}: frecuencias por encima de
$f_s/2$ (frecuencia de Nyquist) no pueden distinguirse de sus imágenes espectrales
más bajas. Como $f_s = 44100$~Hz y las frecuencias DTMF son como máximo
$1477$~Hz~$\ll f_s/2 = 22050$~Hz, este efecto no afecta nuestro análisis.

% ─────────────────────────────────────────────────────────────
\section*{Punto 1 — Identificación de intervalos por dígito}

\subsection*{\texttt{tlfn-a.wav}}

La forma de onda (Figura~\ref{fig:waveform_a}) exhibe bloques de alta amplitud
separados por silencios. Se calculó la energía de la señal en ventanas de 20~ms:
$$
E_k = \sum_{n=kT}^{(k+1)T-1} f[n]^2,
$$
clasificando como \textit{tono activo} toda ventana cuya energía supera el 4\,\%
del máximo. Esto permitió delimitar los 7 dígitos directamente.

\subsection*{\texttt{tlfn-b.wav}}

La señal presenta ruido de fondo continuo a lo largo de toda su duración
(Figura~\ref{fig:waveform_b}), sin pausas visibles entre dígitos. El enfoque de
energía sólo detectó un único segmento. En su lugar se empleó una
\textbf{ventana deslizante de 60~ms} con paso de 20~ms: cada ventana se analiza
espectralmente de forma independiente y se agrupan las ventanas consecutivas que
detecten el mismo par de frecuencias DTMF válidas. \\[1cm]

\begin{figure}[H]
    \centering
    \includegraphics[width=\textwidth]{waveform_tlfn-a.png}
    \caption{Forma de onda de \texttt{tlfn-a.wav}. Los tonos activos están
    sombreados en naranja con el dígito identificado.}
    \label{fig:waveform_a}
\end{figure}

\begin{figure}[H]
    \centering
    \includegraphics[width=\textwidth]{waveform_tlfn-b.png}
    \caption{Forma de onda de \texttt{tlfn-b.wav}. La señal continua con ruido
    requirió ventana deslizante para separar los dígitos.}
    \label{fig:waveform_b}
\end{figure}

% ─────────────────────────────────────────────────────────────
\section*{Punto 2 — Aplicación de la FFT}

Para cada segmento identificado se calcula la DFT usando la FFT. El módulo de los
coeficientes, $|\hat{f}[k]|$, indica la energía presente en la frecuencia
$\xi_k = k \cdot f_s / N$. La búsqueda se restringe al rango $620$--$1550$~Hz,
que cubre todas las frecuencias DTMF y descarta componentes de baja frecuencia
e imagen.

Se retienen los dos índices $k_1, k_2$ con mayor $|\hat{f}[k]|$ en ese rango.
Las frecuencias correspondientes son
$$
\xi_{k_1} = \frac{k_1 \cdot f_s}{N}, \qquad \xi_{k_2} = \frac{k_2 \cdot f_s}{N}.
$$

Nótese que la resolución finita \eqref{eq:resolucion} implica que los picos no
caen exactamente en los valores teóricos de la tabla DTMF, sino a distancias del
orden de $\Delta\xi = f_s/N$. Por eso el mapeo final se realiza buscando el valor
canónico más cercano, con una tolerancia de $\pm 80$~Hz.

La identidad de Plancherel (Teorema~14.6.2) garantiza que esta representación
espectral preserva la energía de la señal (salvo la constante $1/N$):
$$
\sum_{k=0}^{N-1} |f[k]|^2 = \frac{1}{N} \sum_{k=0}^{N-1} |\hat{f}[k]|^2.
$$
Esto asegura que los picos dominantes del espectro corresponden efectivamente a
los componentes de mayor energía en la señal original.

\begin{figure}[H]
    \centering
    \includegraphics[width=\textwidth]{fft_tlfn-a.png}
    \caption{$|\hat{f}[k]|$ para cada dígito de \texttt{tlfn-a.wav}.
    Las líneas punteadas indican los dos picos dominantes y sus frecuencias.}
    \label{fig:fft_a}
\end{figure}

\begin{figure}[H]
    \centering
    \includegraphics[width=\textwidth]{fft_tlfn-b.png}
    \caption{$|\hat{f}[k]|$ para cada dígito de \texttt{tlfn-b.wav}.
    Las líneas punteadas indican los dos picos dominantes y sus frecuencias}
    \label{fig:fft_b}
\end{figure}

% ─────────────────────────────────────────────────────────────
\section*{Punto 3 — Reconocimiento del número marcado}

\subsection*{\texttt{tlfn-a.wav}}

\begin{table}[H]
\centering
\begin{tabular}{ccccc}
\toprule
\textbf{Tono} & \textbf{$\xi_{k_1}$ (Hz)} & \textbf{$\xi_{k_2}$ (Hz)} &
\textbf{DTMF canónico} & \textbf{Dígito} \\
\midrule
1 & 700{,}0  & 1337{,}5 & 697 / 1336 & \textbf{2} \\
2 & 770{,}0  & 1480{,}0 & 770 / 1477 & \textbf{6} \\
3 & 770{,}0  & 1340{,}0 & 770 / 1336 & \textbf{5} \\
4 & 700{,}0  & 1475{,}0 & 697 / 1477 & \textbf{3} \\
5 & 850{,}0  & 1337{,}5 & 852 / 1336 & \textbf{8} \\
6 & 775{,}0  & 1337{,}5 & 770 / 1336 & \textbf{5} \\
7 & 770{,}0  & 1210{,}0 & 770 / 1209 & \textbf{4} \\
\bottomrule
\end{tabular}
\caption{Frecuencias detectadas y mapeo DTMF para \texttt{tlfn-a.wav}.}
\end{table}

\[
\boxed{\texttt{tlfn-a} \;\rightarrow\; \mathbf{2653854}}
\]

\subsection*{\texttt{tlfn-b.wav}}

La ventana deslizante identifica grupos de ventanas con respuesta DTMF consistente,
descartando las regiones donde no se encuentran dos picos claros en la banda de interés.

\begin{table}[H]
\centering
\begin{tabular}{cccc}
\toprule
\textbf{Tono} & \textbf{Intervalo (s)} & \textbf{DTMF canónico} & \textbf{Dígito} \\
\midrule
1 & 0{,}02 -- 0{,}32 & 852 / 1209 & \textbf{9} \\
2 & 0{,}34 -- 0{,}64 & 852 / 1209 & \textbf{7} \\
3 & 0{,}96 -- 1{,}28 & 852 / 1209 & \textbf{7} \\
4 & 1{,}28 -- 1{,}60 & 770 / 1477 & \textbf{6} \\
5 & 2{,}22 -- 2{,}54 & 852 / 1336 & \textbf{8} \\
6 & 2{,}88 -- 3{,}18 & 770 / 1209 & \textbf{4} \\
7 & 3{,}50 -- 3{,}82 & 941 / 1336 & \textbf{0} \\
8 & 3{,}82 -- 4{,}16 & 697 / 1477 & \textbf{3} \\
\bottomrule
\end{tabular}
\caption{Detección por ventana deslizante para \texttt{tlfn-b.wav}.}
\end{table}

\[
\boxed{\texttt{tlfn-b} \;\rightarrow\; \mathbf{97768403}}
\]

% ─────────────────────────────────────────────────────────────
\section*{Conclusiones}

La DFT, calculada eficientemente mediante la FFT ($O(N\log N)$ en lugar de $O(N^2)$),
permite descomponer cada segmento de audio en sus frecuencias constituyentes y
así identificar el par de frecuencias DTMF que lo codifica.

Las diferencias entre los picos detectados y los valores teóricos (del orden de
2--5~Hz) son una consecuencia directa de la resolución espectral finita
$\Delta\xi = f_s/N$: los coeficientes de la DFT sólo toman valores en la grilla
discreta de frecuencias, y no en los valores exactos de la tabla. Segmentos más
largos reducen $\Delta\xi$ y acercan los picos a los valores canónicos.

La elección del método de segmentación resultó determinante. Para
\texttt{tlfn-a.wav}, con silencios claros entre dígitos, fue suficiente un
umbral de energía. Para \texttt{tlfn-b.wav}, cuya señal es continua con ruido
de fondo, fue necesario recurrir a la ventana deslizante, que evalúa la
estructura espectral localmente y desecha las ventanas donde no se detectan dos
picos DTMF válidos.

Los números identificados fueron:
\[
\texttt{tlfn-a} \;\rightarrow\; \mathbf{2653854} \qquad
\texttt{tlfn-b} \;\rightarrow\; \mathbf{97768403}
\]

% ─────────────────────────────────────────────────────────────
\appendix
\section*{Apéndice — Código Python (fragmento clave)}

\begin{lstlisting}[caption={DFT vía FFT de scipy y detección de picos.}]
from scipy.fft import rfft, rfftfreq
from scipy.signal import find_peaks

def detectar_digito_en_segmento(segmento, fs):
    N     = len(segmento)
    # Calcula |f_hat[k]| usando la FFT de scipy (solo frec. positivas)
    fhat  = np.abs(rfft(segmento))
    # Eje de frecuencias: xi_k = k * fs / N
    freqs = rfftfreq(N, d=1/fs)

    # Restringir al rango DTMF
    mask     = (freqs >= 620) & (freqs <= 1550)
    fhat_rec = fhat.copy(); fhat_rec[~mask] = 0

    # Dos picos de mayor |f_hat[k]| -> dos frecuencias dominantes
    picos, _ = find_peaks(fhat_rec,
                          height=fhat_rec.max() * 0.25,
                          distance=int(50 * N / fs))
    top2  = picos[np.argsort(fhat_rec[picos])[::-1]][:2]
    f1, f2 = freqs[top2[0]], freqs[top2[1]]

    dig = digito_dtmf(f1, f2)   # mapeo a tabla DTMF
    return f1, f2, dig
\end{lstlisting}

\end{document}
